% 原书分章 PDF 第 8—10 页：第 28.4 节导言及第 28.4.1—28.4.3 节。
\section{自助法}

在介绍对自助法的推广之前，先简要回顾其理论框架。自助法的基本思想如下：对于估计量 $\hat\Psi:\mathcal M_{NP}\to\boldsymbol\Psi$，我们希望估计其抽样分布
$G(x)=P(\hat\Psi(P_n)\leq x)$；这里，$\hat\Psi(P_n)$ 通过随机的 $P_n$ 被视为随机变量。如果估计量渐近线性，则可用正态分布逼近这一抽样分布，只需估计其一阶矩与二阶矩即可。直接估计这一分布很困难，因为我们只观测到来自 $P_0$ 的一个样本，因而只有 $\psi_n$ 的一个实现值。然而，我们可以从 $P_0$ 的某个估计中，重复抽取 $B$ 个大小为 $n$ 的样本，并对这些样本应用估计量。将这样的样本记为 $O_1^\#,\ldots,O_n^\#$，其对应的经验分布记为 $P_n^\#$，估计记为 $\psi_n^\#=\Psi(P_n^\#)$，则可得到所需抽样分布的估计：
\[
 \hat G(x)=\frac1B\sum_{i=1}^BI(\psi_{n,i}^\#\leq x).
\]

\subsection{非参数自助法}

非参数自助法（nonparametric bootstrap）通过从经验分布 $P_n$ 抽样来应用上述方法。大量研究表明，在各种情形下，它都是估计抽样分布的有效工具。然而，非参数自助法并非在所有情形下都适合估计抽样分布。由于 $P_n$ 是离散分布，从中重复抽样会产生“复制”的观测，即一个自助样本中会包含多个完全相同的观测。Bickel 等（1997a）指出，如果估计量对数据集中的重复取值（ties）敏感，自助法就可能失效。利用交叉验证来权衡偏差与方差的估计量，就是可能对重复取值敏感的一类估计量。如果对非参数自助样本应用交叉验证，重复观测可能同时出现在某一划分的训练部分与测试部分中。这会给依赖交叉验证的估计量带来问题。Hall（1992）专门指出了基于交叉验证的模型选择中的重复取值问题。

这一问题的严重程度取决于预期会出现多少复制观测。对于大小为 $n$、验证比例为 $p_{B_n}$ 的自助样本，一个验证观测在训练样本中具有副本的概率为
\[
 p(\mathrm{copy})=1-\left(1-\frac1n\right)^{(1-p_{B_n})n}
 \approx1-e^{-(1-p_{B_n})}.
\]
十折交叉验证中，$p_{B_n}=0.1$，因此预期约有 $59\%$ 的验证观测也会出现在训练样本中。于是，自助样本上的交叉验证风险估计，平均有 $59\%$ 的部分近似于回代风险估计，从而具有次优性质。对于重复观测给交叉验证带来的问题，一种临时解决方法是采用“按簇”交叉验证（clustered cross-validation）：将完全相同的自助观测组成一簇，再把这些簇分配到训练集与验证集中。这样，就没有观测会同时出现在训练集和测试集中。尽管这一方法缺乏严格论证，我们仍在下面的模拟研究中对其进行了评价。结果中将它标为“NP Bootstrap + CVID”，即“非参数自助法 + CVID”。

\subsection{基于模型的自助法}

参数自助法从基于某个假定参数模型的 $P_0$ 估计 $P_{n,\beta}$ 中抽样。参数自助法可以推广为“基于模型”的自助法（model-based bootstrap），对似然各因子使用半参数或非参数估计。在特定处理下均值的例子中，这意味着使用 $P(Y\mid A,W)$ 和 $\bar g$ 的估计。如果 $Y$ 为二元变量，如本章模拟中所设定的那样，则
$P(Y=1\mid A,W)=E(Y\mid A,W)=\bar Q(A,W)$。
否则，我们需要估计给定 $A,W$ 时 $\epsilon(A,W)$ 的条件分布，其中
$Y=E(Y\mid A,W)+\epsilon(A,W)$。
具体地，按下列算法从估计 $\tilde P_n=Q_{W,n}g_nQ_{Y,n}$ 中抽取观测。

\noindent\begin{minipage}{\linewidth}
\begin{sourcebox}[算法。基于模型的自助法]
\begin{itemize}[label=$\star$,leftmargin=1.6em]
\item 从 $W$ 的经验分布 $Q_{W,n}$ 中抽取 $W^\#$；
\item 利用 $W^\#$，从 $g_n(\cdot\mid W^\#)$ 中抽取 $A^\#$；
\item 利用 $A^\#$ 与 $W^\#$，从
$Q_{Y,n}(\cdot\mid(A^\#,W^\#))$ 中抽取 $Y^\#$。
\end{itemize}
\end{sourcebox}
\end{minipage}\par

下一节介绍的靶向自助法，是一种特定的基于模型的自助法，它使用经靶向更新的估计 $\bar Q_n^*$、$\bar g_n^*$，以确保正确的渐近表现。

\subsection{靶向自助法}

靶向自助法的思想是构造 TMLE $\tilde P_n^*$，使
$\Sigma^2(\tilde P_n^*)$ 成为 $\sigma_0^2=\Sigma^2(P_0)$ 的 TMLE。
于是我们知道，在正则条件下，$\Sigma^2(\tilde P_n^*)$ 是 $\sigma_0^2$ 在 $P_0$ 处的渐近线性且有效的估计量，因此可以构造 $\sigma_0^2$ 的置信区间。此外，由于它是 $\sigma_0^2$ 的代入估计量，在有限样本中通常更加可靠，因而可能改善置信区间的有限样本覆盖率。另外，我们将说明，在适当正则条件下，由于 $\Sigma^2(\tilde P_n^*)$ 的相合性，基于抽样
$O_1^\#,\ldots,O_n^\#\sim_{iid}\tilde P_n^*$ 得到的
\[
 \sqrt n\bigl(\hat\Psi(P_{n,\#})-\Psi(\tilde P_n^*)\bigr)
\]
的自助分布，在给定几乎每一个 $(P_n:n)$ 时，
$\sqrt n(\hat\Psi(P_n)-\psi_0)$ 收敛到所需的极限分布
$N(0,\sigma_0^2)$，即使 $\tilde P_n^*$ 设定错误也是如此。因此，可以证明，为估计 $\sqrt n(\psi_n-\psi_0)$ 的极限分布，
$\tilde P_n^*$ 自助法几乎处处相合；自助法还有一个优点，即也能估计该估计量在这一自助分布下的有限样本抽样分布。通常，从 $P_0$ 的估计 $\tilde P_n^*$ 中抽样的基于模型的自助法，其相合性依赖于 $\tilde P_n^*$ 的密度作为 $P_0$ 密度估计量的相合性。然而在这里，自助法的相合性仅依赖于使 TMLE $\Sigma^2(\tilde P_n^*)$ 成为 $\Sigma^2(P_0)$ 的相合估计量的条件。因此，在 $\tilde P_n^*$ 中，可以允许对 $P_0$ 的某些部分作出不相合的估计。我们将这一自助法称为靶向自助法。

$\Sigma^2(P_0)$ 的 TMLE 通常可表示为 $\Sigma^2(Q_{\Sigma,n}^*)$，其中采用了较小的参数 $P\mapsto Q_\Sigma(P)$，并可能用到 $g_\Sigma(P_0)$ 的干扰参数估计量 $g_{\Sigma,n}^*$。
因此，$\tilde P_n^*$ 可以定义为任意满足
$Q_\Sigma(\tilde P_n^*)=Q_n^*$ 的分布，而不影响靶向自助法的相合性。这表明靶向自助法对某些类型的模型设定错误具有稳健性。为在估计 $\hat\Psi(P_n)$ 的实际抽样分布时取得最佳的有限样本表现，除了 $Q_0$ 之外，还能很好地逼近 $P_0$ 的其余部分，可能仍有帮助；不过，这一贡献在渐近意义下可以忽略。
